\section{Binary transport and bit bounds for sector transcripts}
\label{sec:sector-integer-transport}

The compressed hierarchy needs to operate on binary normal coordinates,
not on an expanded list of normal discs. A local verifier must also accept
its own exact serialization. The retained sector-exhaustion checker had a
mismatch: its disc subproofs accepted signed hexadecimal integer fields,
but its ray coordinates and Euler values required Python integers. Thus an
otherwise valid serialized transcript could fail without any mathematical
change. The nested Q-support reader also used the truth value of a field,
rather than its integer value: a string representing zero is nonempty.

The repair decodes only the ray-count and Euler fields with the existing
strict integer codec. Integers and explicitly signed hexadecimal strings
are accepted; booleans, floats, decimal strings and malformed encodings
are rejected. Coordinates must still have the independently computed
support width. Decoded rays are compared with the independently rebuilt
complete ray set, with the same duplicate, Euler and component-proof
checks. The already verified nested Q transcript is decoded before its
positive support is extracted. Source identifiers, allowed type indices,
phase and status remain subject to their existing schema checks.
No producer, claimed complete enumeration or claimed count is trusted.

\subsection{The size of primitive ray fields}

The following determinant argument is the usual vertex-coordinate bound;
we include it to make the implementation's bit accounting explicit.
Lackenby and Yazdi recall the $2^{7t-1}$ bound and attribute it to Hass
and Lagarias\cite[Section 2.7]{lackenby-yazdi-genus}.
This is not a new knot-theory theorem.

\begin{proposition}[Primitive sector ray size]
For a source triangulation with $t$ tetrahedra, a primitive integral
extreme ray of its standard matching cone restricted to an admissible
support has every coordinate at most $2^{7t-1}$. Its complete coordinate
vector therefore has $O(t^2)$ bits, and its Euler characteristic has
$O(t+\log(t+1))$ bits. The ray-coordinate and Euler fields of a transcript
with $R$ such rays have $O(Rt^2)$ bits.
\end{proposition}
\begin{proof}
Each original source matching row equates one triangle and one
quadrilateral contribution on the two sides of a face. Before cancellation
it has two positive and two negative unit terms. The positive triangle and
quadrilateral indices are distinct, as are the two negative indices.
Identifications can cancel opposite terms but cannot combine equal-sign
terms. Thus each row has Euclidean norm at most two. Deleting columns
preserves this bound.

Let the ray have $s$ positive coordinates, with $s\le7t$. The matching
matrix restricted to those columns has rank $s-1$: otherwise its
nonnegative cone contains a perturbation of the positive ray in another
direction, contradicting extremality. Select $s-1$ independent rows.
The vector of signed maximal minors spans their one-dimensional kernel.
Hadamard's inequality bounds each minor by $2^{s-1}$. Dividing by their
common divisor and choosing the nonnegative sign gives the primitive
ray and preserves this bound. For $s=1$ the primitive entry is one.

There are at most $7t\,2^{7t-1}$ normal discs in the resulting surface.
Each has at most four sides and corners. Counting cells gives, for example,
$|\chi|\le9\cdot7t\,2^{7t-1}$, sufficient for the stated digit bound.
The vector has at most $7t$ entries, each with at most $7t$ binary digits;
its quadrilateral subvector and Euler field satisfy the same accounting.
\end{proof}

The bound applies to primitive rays, not arbitrary multiples or arbitrary
normal surfaces. The dense reference checker confirms a candidate's
one-dimensional supported kernel and compares its primitive direction.
The implementation never expands its represented normal discs. Signed
hexadecimal transport uses a constant number of characters per four bits,
so decoding these fields adds time linear in their encoded size to the
existing exact arithmetic work. A maliciously changed 16,385-bit count
must still fail the geometric ray equality; a large number is not evidence
of a new ray.

This is not a polynomial bound on the entire exhaustion certificate.
The number $R$ of rays can be large; negative positive-Euler ray entries
also carry independently replayed component proofs, whose size is accounted
separately. The checker still reconstructs the reference ray set and can
perform a large arrangement enumeration. Small binary coordinate fields
alone bound neither search coverage nor hierarchy depth, source growth,
restarts, or component-proof generation. In particular no general
subexponential recognition bound follows from this repair.

\subsection{Real-source replay evidence}
All 1,534 maintained tests pass in 293.776 seconds; the 24 focused tests
pass in 2.474 seconds. Real-source regressions cover all three exhaustion
statuses, equivalent signed hexadecimal values, mixed representations,
nested Q support containing encoded zeros, and producer-disabled replay.
Malformed numbers, altered large counts, duplicate rays, incorrect widths,
changed Euler values and encoded metadata indices are rejected. Cancellation
continues to propagate.

The retained-source audit replays 134 saved exhaustion certificates:
twelve generic-Q exclusions, 61 feasible-span proofs and 61 matching-support
proofs. These are proof instances, not 134 distinct triangulations.
Every literal certificate passes both the frozen preceding checker and the
new checker. Every equivalent hexadecimal transcript fails the frozen
checker and passes the new one; every signed-zero variant also passes the
new checker. All discovery, kernel, ray and disc-count producers are disabled
during replay. The existing dense geometric reconstruction and independent
negative component proof checks remain active.

For every retained proof the audit rejects an altered 8,193-bit Euler
value and an altered 16,385-bit quadrilateral count after ordinary JSON
transport. The latter uses the public binary codec, rather than changing
Python's decimal conversion limit. These changed values are mutations,
not valid high-coordinate geometric witnesses. The audit checks the
primitive quadrilateral digit bound on all retained entries. It does not
claim that this finite bank tests every large triangulation or measures a
new asymptotic speedup.

Runtime release \code{b361cd6cd} changes only the existing independent
sector checker and adds five regression tests. The audit pins 422 current
core-runtime/test sources and checks the preceding 652-source inventory: exactly
the verifier changed among those existing sources. Discovery producers,
recognition policy and disabled-stage behavior are unchanged.

The reproduction driver is \code{data/sector\_integer\_audit.py}.

The records are under \code{data/sector-integer-*}. The publication review checks exact
source hashes, the passing test records and all 134 transport outcomes.
The current \code{openai/math} head is still \code{fd4aeeb2ee}, checked
against its remote during this continuation; no new result is inferred
from that catalogue. The general quasi-polynomial search obligations remain
open.

